\subsection*{Drawing a recorded successor from visible parts}
Removing a preference shared across contexts changed the native writer's
strokes. On the same 512 article pairs, we compared full writing-code density
with its log contrast against the last observed glyph, at the same plausible
component means. The average difference between the two contexts grew from
7.08 to 60.18 pixels per 1,024-pixel drawing. Neither rule reproduced any of
the 1,024 successor images. Mean source ink F1 fell from .315 to .237.
Both complete outputs about the history of the character for speech remained
unreadable.

A separate model generated sixteen physical patches per glyph. From the two
most recent visible images, it estimated how likely each patch was to appear
at its position and to join the patches already drawn. On 668 new contexts
from 334 further articles, the model trained on authentic passage order drew
18 exact successors spanning 14 characters (Fig.~\ref{fig:patch-writing}).
Moving each training successor to another position in the same passage gave
zero exact successors. The difference was 2.69 percentage points (95\%
paired article interval 1.50 to 3.89). The model drew 198 images identical
to known nonempty glyphs and left 348 outputs blank. Neither model got both
successors right in any article. The difference establishes that local order
can help a writer produce the source image without seeing it among candidate
answers.

Fitting the ordered count tables after preparation took 0.75 seconds on CPU.
Both fits, preparation, checkpoint writing, reload and evaluation took 23.14
seconds. The ordered checkpoint occupied 18.98 MB compressed. The three free
article continuations and the answer about the history of the character for
speech reached 128 cells and fell into repetition. Accurate occasional
successors leave sustained writing and useful answer generation open.
